\section*{Notation and scope}
\addcontentsline{toc}{section}{Notation and scope}
Our exit arrows move from a deeper stratum to a shallower one.  All homotopy
groups in Parts 1--5 are discrete unless a profinite completion is explicitly
displayed.  A base point in a stack means a geometric point together with the
usual point of its classifying homotopy type.  We never replace that homotopy
type by the coarse space.

\begin{center}\small
\begin{tabular}{p{.18\textwidth}p{.41\textwidth}p{.32\textwidth}}
\toprule
Our notation & Meaning & Manuscript notation and scan \\
\midrule
$D=\Theta$, $U=X\setminus D$ & Reduced boundary and open complement &
 $\Theta$, $X^*$; 147/106, 127 \\
$D_d$, $D_d^\circ$ & Normalized $d$-branch intersections, and their open part &
 $D_d$, $D_d^*$; 147/106--108, 123, 127 \\
$D_{d_0,\ldots,d_r}^\circ$ & Finite incidence cover of the deepest open stratum &
 $D_{d_*}^*$; 147/128, (205)--(206) \\
$\mathcal I$, $\mathcal I_\bullet$ & Internal category of branch flags, and its nerve &
 $\mathcal D$, $\mathcal D_*$; 147/114, 128 \\
$\Lambda_s$ & Integral lattice of local meridians &
 $T_d$; 147/121--123, (183), (192) \\
$P_s$, $L_s$ & Punctured multinormal bundle; its unit-torus bundle &
 $P_d^*$; 147/122--123, (189)--(191) \\
$L_{ST}$ & Relative torus link over the incidence cover for $S<T$ &
 $\breve N^{**}$; 147/177--178, (295) \\
$M=X^{\log}$, $L=\partial M$ & Real oriented blow-up and divisor link &
 No name for $M$; $X^{*\infty}$ for the link germ, 147/150 \\
$\mathcal E=\Exit(X,D)$ & Exit-path $\infty$-category &
 No identical manuscript symbol; reconstructed from flags and tubes \\
$h\mathcal E$, $B\mathcal E$ & Ordinary homotopy category; groupoidification &
 Distinction required by 147/179--180 \\
$\mathsf T$ & Coherent diagram of elementary tube homotopy types &
 $\mathcal V^{d'_0}_{d_*,d''_0}$; 147/168, 177--178 \\
$\Delta(D)$ & Coarse dual complex, using only nonempty branch sets &
 ``maquette simpliciale''; 149/133 \\
\bottomrule
\end{tabular}
\end{center}

The table fixes the direction of every flag: $d_0>\cdots>d_r$ means
deepest to shallowest.  It avoids the exchange of $d_*$ and $d'_*$ explicitly
announced in the margin of 147/135.  ``Canonical'' means a natural object or
map in the appropriate homotopy category or $\infty$-category, with the stated
coherences.  A path from a local base point to a global one changes the induced
homomorphism by conjugation.  We do not call all natural isomorphisms between
two such functors unique: even the identity of $B\Z$ has automorphism group
$\Z$.

We use paracompact complex manifolds and smooth analytic Deligne--Mumford
stacks presented by an \emph{étale} analytic groupoid of such manifolds.
Normal crossings are checked in this presentation; branches may be permuted,
and global components may have self-intersections.  Classifying spaces and
constructible sheaves on the stack are obtained by descent along the whole
Čech nerve.  No properness of the diagonal is used in the local torus
calculation.  When a result needs a separated finite-isotropy orbifold, a
proper completion, an algebraic variety, or a mixed Hodge structure, we state
that extra hypothesis.  In particular, the infinite-$\Z$ transvection quotient
in 147/146 is not a separated finite-isotropy orbifold.

The empty divisor is allowed.  Assertions concerning a nonempty link assume
$D\ne\varnothing$.  A link for a nonproper pair is the germ of a punctured
neighbourhood of this specified divisor, not a claim about every end of $U$.
All constructions are componentwise.  There are no exit arrows between
unrelated strata; an incidence cover can have several components, each of
which must be retained.

The special curve types are not suppressed.  A sphere with zero punctures
is not aspherical; with one puncture it has trivial peripheral monodromy;
with two punctures its meridians are infinite; a compact elliptic curve is
aspherical with no boundary test.  The stable moduli notation is used only
when $2g-2+n>0$: $(0,3)$ is a point and $(1,1)$ is a stack, while $(0,0)$,
$(0,1)$, $(0,2)$, and the unpointed genus-one moduli problem do not define the
stable Deligne--Mumford pair in question.  Real and nonorientable pairs are
outside the stated complex hypotheses; their normal links need not be
circles.  The arithmetic results below are in characteristic zero and make
no claim about wild reduction in characteristics $2$ or $3$.

We have checked the supplied STRATA claims against their proofs where used,
and have opened the external results cited below.  Neither the dictionary
nor a companion AI return is used as an authority.  Numbered statements carry
a status; ``Known'' is used only when the displayed reduction to an opened
external theorem is supplied.  This work is AI-produced by GPT-6 Astra
(OpenAI, ChatGPT Pro mode), attributes no authorship to a person, and makes
no priority claim.
